\documentclass[10pt,a4paper]{amsart}
\usepackage{amsmath,amssymb,amsthm,graphicx}
\usepackage{verbatim}
\newenvironment{reference}{\comment}{\endcomment}
%\newenvironment{reference}{}{}
\newenvironment{slogan}{\comment}{\endcomment}
\newenvironment{history}{\comment}{\endcomment}

% For commutative diagrams we use Xy-pic
\usepackage[all]{xy}

% We use 2cell for 2-commutative diagrams.
\xyoption{2cell}
% Theorem environments.
%
\theoremstyle{plain}
\newtheorem{theorem}[subsection]{Theorem}
\newtheorem{proposition}[subsection]{Proposition}
\newtheorem{lemma}[subsection]{Lemma}

\theoremstyle{definition}
\newtheorem{definition}[subsection]{Definition}
\newtheorem{example}[subsection]{Example}
\newtheorem{exercise}[subsection]{Exercise}
\newtheorem{situation}[subsection]{Situation}

\theoremstyle{remark}
\newtheorem{remark}[subsection]{Remark}
\newtheorem{remarks}[subsection]{Remarks}

\numberwithin{equation}{subsection}

% Macros
%
\def\lim{\mathop{\mathrm{lim}}\nolimits}
\def\colim{\mathop{\mathrm{colim}}\nolimits}
\def\Spec{\mathop{\mathrm{Spec}}}
\def\Hom{\mathop{\mathrm{Hom}}\nolimits}
\def\Ext{\mathop{\mathrm{Ext}}\nolimits}
\def\SheafHom{\mathop{\mathcal{H}\!\mathit{om}}\nolimits}
\def\SheafExt{\mathop{\mathcal{E}\!\mathit{xt}}\nolimits}
\def\Sch{\mathit{Sch}}
\def\Mor{\mathop{\mathrm{Mor}}\nolimits}
\def\Ob{\mathop{\mathrm{Ob}}\nolimits}
\def\Sh{\mathop{\mathit{Sh}}\nolimits}
\def\NL{\mathop{N\!L}\nolimits}
\def\CH{\mathop{\mathrm{CH}}\nolimits}
\def\proetale{{pro\text{-}\acute{e}tale}}
\def\etale{{\acute{e}tale}}
\def\QCoh{\mathit{QCoh}}
\def\Ker{\mathop{\mathrm{Ker}}}
\def\Im{\mathop{\mathrm{Im}}}
\def\Coker{\mathop{\mathrm{Coker}}}
\def\Coim{\mathop{\mathrm{Coim}}}

% Boxtimes
%
\DeclareMathSymbol{\boxtimes}{\mathbin}{AMSa}{"02}

%
% Macros for moduli stacks/spaces
%
\def\QCohstack{\mathcal{QC}\!\mathit{oh}}
\def\Cohstack{\mathcal{C}\!\mathit{oh}}
\def\Spacesstack{\mathcal{S}\!\mathit{paces}}
\def\Quotfunctor{\mathrm{Quot}}
\def\Hilbfunctor{\mathrm{Hilb}}
\def\Curvesstack{\mathcal{C}\!\mathit{urves}}
\def\Polarizedstack{\mathcal{P}\!\mathit{olarized}}
\def\Complexesstack{\mathcal{C}\!\mathit{omplexes}}
% \Pic is the operator that assigns to X its picard group, usage \Pic(X)
% \Picardstack_{X/B} denotes the Picard stack of X over B
% \Picardfunctor_{X/B} denotes the Picard functor of X over B
\def\Pic{\mathop{\mathrm{Pic}}\nolimits}
\def\Picardstack{\mathcal{P}\!\mathit{ic}}
\def\Picardfunctor{\mathrm{Pic}}
\def\Deformationcategory{\mathcal{D}\!\mathit{ef}}
\usepackage{cleveref}
\crefname{originallemma}{original lemma}{original lemmas}
\Crefname{originallemma}{Original lemma}{Original lemmas}
\makeatletter
\newlabel{lemma-additive-additive}{{0105}{source}{}{sourcepointer-1}{}}
\newlabel{lemma-additive-additive@cref}{{[originallemma][0][]0105}{[1][0][]source}{}{}{}}
\makeatother
\title{A functor between additive categories is additive iff it preserves direct sums}
\author{The Stacks Project Authors}
\date{stacks-project:homology.tex:local-snapshot-20261012; GNU FDL 1.2 or later}
\begin{document}
\maketitle
\noindent\textit{Editorial note.} All delivered mathematical text is the author's own, copied from the pinned Stacks Project chapter file homology.tex; no statement or proof step was rewritten, reordered or removed.\par\medskip
\noindent\textit{Editorial note.} The proof is the author's complete argument for all three conditions; the author's introductory sentence describing the lemma as 'completely silly' sits outside the counted spans, and every step of the argument inside them is written out and finishes with a conclusion rather than a placeholder.\par\medskip
\noindent\textit{Editorial note.} Every labelled reference occurring inside the counted spans is carried in assembly.external\_references as a non-load-bearing navigation pointer (same-chapter ones at their definition site, cross-chapter ones at the invocation inside the counted span), so no load-bearing mathematics is replaced by a stub. No bibliography citation occurs inside the counted spans.\par\medskip
\noindent\textit{Editorial note.} The same-name check for this card scanned every staged Stacks chapter of this repository -- all sources/stacks-*/pinned-original-source.tex files, the card's own chapter excluded -- and the labels found are recorded in dedup.same\_name\_cross\_chapter on the card.\par\medskip
\noindent\textit{Editorial note.} The preamble is the chapter's own native preamble excerpt already used for this source: the comment-type reference, slogan and history environments (native preamble lines 8-13), the Xy-pic setup (native preamble lines 14-18) and the native theorem declarations with the mathematical macros (native preamble lines 154-217, the actual last line of the file).\par\medskip
\noindent\textit{Allowed standard prerequisites (editorial).}\begin{itemize}
\item Additive functors commute with direct sums. Same chapter, Lemma lemma-additive-additive. Gives the implication from additivity to the preservation of binary direct sums in either direction.
\item Addition of morphisms through the diagonal and the codiagonal. Same chapter, the description of f + g as the composite A -> A + A -> B + B -> B via the diagonal and the codiagonal. Supplies the description of the sum of two morphisms to which the functor is applied in the final step of the proof.
\end{itemize}\medskip
\ifdefined\hypertarget\hypertarget{sourcepointer-1}{}\fi
\noindent\textit{Source reference.} Original lemma 0105: sources/stacks-homology/pinned-original-source.tex, source lines 390--390. Navigation pointer invoked at proof line 1702; the counted span contains the pointer only and the cited result is not reproduced inside it, so the pointer is registered as a non-load-bearing navigation stub in the established form for Stacks material. No bibliography citation is involved.\par\medskip
\begin{lemma}
\label{lemma-additive-functor}
Let $\mathcal{A}$ and $\mathcal{B}$ be additive categories.
Let $F : \mathcal{A} \to \mathcal{B}$ be a functor.
The following are equivalent
\begin{enumerate}
\item $F$ is additive,
\item $F(A) \oplus F(B) \to F(A \oplus B)$ is an isomorphism for
all $A, B \in \mathcal{A}$, and
\item $F(A \oplus B) \to F(A) \oplus F(B)$  is an isomorphism for
all $A, B \in \mathcal{A}$.
\end{enumerate}
\end{lemma}

\begin{proof}
Additive functors commute with direct sums by
Lemma \ref{lemma-additive-additive} hence (1)
implies (2) and (3). If (2) holds and $0$ denotes the
zero object of $\mathcal{A}$, then the map
$F(0) \oplus F(0) \to F(0 \oplus 0) \cong F(0)$
being an isomorphism implies that the diagonal map
$\Hom(F(0), C) \to \Hom(F(0), C) \oplus \Hom(F(0), C)$
is bijective for all $C$ in $\mathcal{B}$. Thus $F(0)$
is the zero object of $\mathcal{B}$ and $F$ maps the
zero morphism between any two objects to the zero morphism.
Similarly holds if (3) holds. It follows that (2) and (3)
are equivalent because the composition
$F(A) \oplus F(B) \to F(A \oplus B) \to F(A) \oplus F(B)$
is given by the matrix
$$
\left(
\begin{matrix}
F(1) & F(0) \\
F(0) & F(1)
\end{matrix}
\right)
$$
and therefore is the identity map. Assume (2) and (3) hold.
Let $f, g : A \to B$ be maps. Then $f + g$ is equal to
the composition
$$
A \to A \oplus A \xrightarrow{\text{diag}(f, g)} B \oplus B \to B
$$
Apply the functor $F$ and consider the following diagram
$$
\xymatrix{
F(A) \ar[r] \ar[rd] &
F(A \oplus A) \ar[rr]_{F(\text{diag}(f, g))} & &
F(B \oplus B) \ar[r] \ar[d] &
F(B) \\
&
F(A) \oplus F(A) \ar[u] \ar[rr]^{\text{diag}(F(f), F(g))} & &
F(B) \oplus F(B) \ar[ru]
}
$$
We claim this is commutative. For the middle square we can verify it
separately for each of the four induced maps $F(A) \to F(B)$
where it follows from the fact that $F$ is a functor and
transforms zero morphisms into zero morphisms.
For the triangle on the left, we use that
$F(A \oplus A) \to F(A) \oplus F(A)$ is an isomorphism
to see that it suffice to check after composition with
this map and this check is trivial. Dually for the other triangle.
Thus going around the bottom is equal to $F(f + g)$ and we conclude.
\end{proof}

\end{document}
